\section{Exact relations}\label{sec:8.3}

A computed response of silicon can be wrong in ways that no convergence test reveals. Take an electron--nucleus interference spectrum that keeps one sign, an energy-weighted spectrum that counts thirteen electrons per atom, or a coupling proportional to mass that excites optical phonons at vanishing momentum transfer. Each is an error, whatever the method and however well converged, because each violates a relation that holds for every target. \Cref{sec:4.2,sec:4.3} derived the first of these relations: detailed balance, the positivity of the response matrix, the fluctuation--dissipation theorem, the Kramers--Kronig relations and the f-sum rule of every species. The multi-operator rate of \cref{sec:8.2} needs three more: an interference bound, higher moments and a Ward identity. Positivity limits how strongly two couplings can interfere. The higher frequency moments of the response matrix are fixed by commutators, and the third moment of the electron--nucleus interference turns out to be set by the electron density at the nuclei. Translation invariance links the densities of all species to the momentum density, and as a result a coupling proportional to mass can only make sound at long wavelengths. None of these relations depends on an approximation to the target, so a response that violates one is wrong, not merely truncated.

\subsection{Interference of two couplings}\label{sec:8.3-bound}

In the polarized electron gas of \cref{we:8-cross} the rate of a density and a spin coupling together was $S_\uparrow(\tilde V_n + \frac12\tilde V_s)^2$, which can be anything from zero to twice the sum of the two single-operator rates. How far can a cross term go in general? Split the rows of \cref{eq:8.2-rate} into two groups, $A$ and $B$, for instance the electron rows and the nuclear rows of a dark photon, or the density and the current rows of \cref{sec:8.1}. The rate is then
\begin{equation*}
  \varGamma = \varGamma_A + \varGamma_{AB} + \varGamma_B ,
\end{equation*}
where $\varGamma_A$ and $\varGamma_B$ keep only the terms with both indices in one group, the rates of each group alone, and the cross term $\varGamma_{AB}$ collects the terms with one index in each, $2\Real\sum_{a\in A,b\in B}\mathfrak{D}_{ab}S_{ab}$ under the integral.

The bound follows from a family of models. Multiply every coupling of group $B$ by a real number $\lambda$. Each element $\mathfrak{D}_{ab}$ is multiplied by $\lambda$ once for every index in $B$, which is $\mathfrak{D}$ multiplied on both sides by a real diagonal matrix, so the tensor is still Hermitian and positive semidefinite. Its contraction with the positive semidefinite $S_{ab}$ of \cref{eq:4.2-positivity} is therefore never negative, at every $(\vect{q},\omega)$ and for every velocity. So is the rate after any integration over $\vect{q}$, $\omega$ and the halo:
\begin{equation*}
  \varGamma(\lambda) = \varGamma_A + \lambda\,\varGamma_{AB} + \lambda^2\,\varGamma_B \ge 0 \qquad \text{for every real } \lambda .
\end{equation*}
A quadratic that never goes negative has a discriminant that is not positive, $\varGamma_{AB}^2 \le 4\varGamma_A\varGamma_B$.
\begin{keyresult}[Interference bound]
For any two groups of rows, any target and any dark-matter model,
\begin{equation}
  |\varGamma_{AB}| \le 2\,(\varGamma_A\varGamma_B)^{1/2} , \qquad\text{equivalently}\qquad
  \bigl(\varGamma_A^{1/2} - \varGamma_B^{1/2}\bigr)^2 \le \varGamma \le \bigl(\varGamma_A^{1/2} + \varGamma_B^{1/2}\bigr)^2 .
  \label{eq:8.3-bound}
\end{equation}
The cross term is at most twice the geometric mean of the separate rates. The bound holds pointwise in $(\vect{q},\omega)$, for the spectrum $\dd\mathcal{R}/\dd\omega$ in any bin, and for the total rate.
\end{keyresult}
\noindent For one row in each group the pointwise statement is the Cauchy--Schwarz inequality $|S_{ab}|^2 \le S_{aa}S_{bb}$ of \cref{eq:4.2-positivity} combined with $|\mathfrak{D}_{ab}|^2 \le \mathfrak{D}_{aa}\mathfrak{D}_{bb}$ from \cref{eq:8.1-D}, and the integrated statement then follows from the Cauchy--Schwarz inequality for integrals, $\int(xy)^{1/2} \le (\int x\int y)^{1/2}$. The argument through $\lambda$ gives both at once, and for blocks of any size.

The bound is reached when both matrices have rank one on the two groups and their phases line up, as in the polarized gas: there $S_{ns}^2 = S_{nn}S_{ss}$, and at $\tilde V_s = -2\tilde V_n$ the rate sits on the lower edge of \cref{eq:8.3-bound}, zero. The bound is not an estimate. For an unpolarized target and couplings that are blind to the dark-matter spin, the cross terms between density and spin vanish by symmetry (\cref{sec:8.2-cross}), far inside the band. Its use is to put a rigorous band on a rate whose interference term has not been computed: writing $r = \varGamma_B/\varGamma_A$, the total lies between $(\sqrt r - 1)^2$ and $(\sqrt r + 1)^2$ times $\varGamma_A$. The same inequality bounds the Migdal interference, $|S_{e\kappa}|^2 \le S_{ee}S_{\kappa\kappa}$, and enters the nuisance parameters of \cref{sec:8.8} as a constraint.

\subsection{Moment matrices}\label{sec:8.3-moments}

The f-sum rule fixed the first frequency moment of every density by a double commutator (\cref{eq:4.3-firstmoment,eq:4.3-species}). Higher moments weight the spectrum toward high energies, and they are where an approximation to the core electrons or to the electron--nucleus interaction shows up. We collect them into matrices,
\begin{equation}
  \mathscr{M}^k_{ab}(\vect{q}) \equiv \int_{-\infty}^{\infty}\frac{\dd\omega}{2\pi}\;\omega^k\bigl[S_{ab}(\vect{q},\omega) + S_{ba}(-\vect{q},\omega)\bigr]
  = \frac{1}{\Vol}\Bigl\langle\bigl[O_a(\vect{q})^\dagger,\,[H, O_b(\vect{q})]_k\bigr]_{\mp}\Bigr\rangle ,
  \label{eq:8.3-moments}
\end{equation}
for operators with $O(-\vect{q}) = O(\vect{q})^\dagger$ and any stationary state, where $[H, O]_k$ is $O$ commuted $k$ times with $H$, as in \cref{sec:4.3}, and $[X, Y]_\mp$ is the commutator $XY - YX$ for odd $k$ and the anticommutator $XY + YX$ for even $k$. The derivation of \cref{eq:4.3-moments} goes through unchanged. The term $O_a^\dagger[H, O_b]_k$ gives the $k$-th moment of $VS_{ab}(\vect{q},\omega)/2\pi$, the term $[H, O_b]_kO_a^\dagger$ gives $(-1)^k$ times that of $VS_{ba}(-\vect{q},\omega)/2\pi$, and the sign of the bracket makes the two add for every $k$. For $k = 1$ this is \cref{eq:4.3-firstmoment}, so $\mathscr{M}^1$ is the first-moment matrix of \cref{sec:4.3}. Since $S_{ab}$ has the units of $\mathrm{eV}^2$ for densities, $\mathscr{M}^k$ has those of $\mathrm{eV}^{k+3}$.

The odd and the even moments carry different information. An odd moment is a commutator. It needs only the algebra of $H$ and equal-time averages of simpler operators, and for the densities of electrons and nuclei the first one is a number, \cref{eq:4.3-species}. An even moment is an equal-time correlation function. On the diagonal, with $X = [H, O_a]_{k/2}$, it reads $\Vol^{-1}\langle X^\dagger X + XX^\dagger\rangle$; for instance $\mathscr{M}^0_{ee} = 2\nd_eS^{\mathrm{st}}$ for a target with inversion symmetry, by \cref{eq:4.2-static}. It depends on the state, its temperature and its correlations, but no dynamics is needed to compute it, which is why the moments of the nuclear self response can be taken from a static path-integral simulation (\cref{app:E.19}). In the same way as for $k = 1$ in \cref{sec:4.3}, an odd $\mathscr{M}^k$ is the coefficient of $z^{-(k+1)}$ in the expansion of $R_{ab}(\vect{q}, z)$ at large $z$ (\cref{app:D.26}).

Positivity constrains the moments among themselves. The spectral weight of one row, summed over $\pm\vect{q}$, is never negative. It defines a positive measure on the real line, and its moments $\mathscr{M}^k_{aa}$ form a Hankel matrix, with elements $\mathscr{M}^{i+j}_{aa}$, that is positive semidefinite:
\begin{equation}
  \sum_{i,j=0}^{n}x_i\,x_j\,\mathscr{M}^{i+j}_{aa} = \int\frac{\dd\omega}{2\pi}\Bigl(\sum_ix_i\omega^i\Bigr)^2\bigl[S_{aa}(\vect{q},\omega) + S_{aa}(-\vect{q},\omega)\bigr] \ge 0 .
  \label{eq:8.3-hankel}
\end{equation}
The smallest case is $(\mathscr{M}^1_{aa})^2 \le \mathscr{M}^0_{aa}\mathscr{M}^2_{aa}$. At $T = 0$ the measure lives on $\omega \ge 0$, and the shifted matrix $\mathscr{M}^{i+j+1}_{aa}$ is positive semidefinite too, which gives for instance $(\mathscr{M}^2_{aa})^2 \le \mathscr{M}^1_{aa}\mathscr{M}^3_{aa}$. A spectrum is determined, within a band that these inequalities make precise, by its first few moments, the Markov--Krein band of \cref{sec:11.5}.

How many moments exist? For electrons and point nuclei with Coulomb interactions, $\mathscr{M}^3$ is finite and $\mathscr{M}^4$ is not. The fourth moment of a density is the equal-time average of $|[H,[H,n_s]]|^2$, and $[H,[H,n_s]]$ contains the force on each particle (\cref{app:D.26}). The mean square of the Coulomb force between an electron and a nucleus, $\int\dd^3y\,\nd(\vect{y})\,(Z_\kappa\alpha/y^2)^2$ with $\nd$ finite at the nucleus, diverges as $\int\dd y/y^2$. The spectrum shows the same thing at large $\omega$. The oscillator strength of the hydrogen continuum falls as $\omega^{-7/2}$ (\cref{we:4-trk}), and at small $q$, where $S \propto q^2\,\omega^{-1}\dd f^{\mathrm{osc}}/\dd\omega$, the structure factor falls as $\omega^{-9/2}$: $\omega^3S$ is integrable and $\omega^4S$ is not. In the Born--Oppenheimer approximation, the order $n_B = 0$ of the expansion in the derivative coupling of \cref{sec:8.5}, the nuclei move on a smooth energy surface, and every moment of their self response is finite.

\subsection{The mixed third moment}\label{sec:8.3-third}

The first moment of the electron--nucleus interference vanishes, \cref{eq:4.3-species}, and at $T = 0$ the real part of $S_{e\kappa}$ must change sign along the energy axis. Which sign wins at high energy, where the Migdal effect and core excitations live? The third moment weights the spectrum by $\omega^3$ and answers the question.

The calculation takes three steps, written out in \cref{app:D.26}. First, in a stationary state the average of $[X, [H, Y]]$ equals that of $[[X, H], Y]$, since the average of a commutator with $H$ vanishes; this moves one commutator onto the electron density, $\mathscr{M}^3_{e\kappa} = \Vol^{-1}\langle[[n_e^\dagger, H], [H, [H, n_\kappa]]]\rangle$. Second, the continuity equations of the two species, \cref{eq:4.3-continuity} and its nuclear analogue, turn the two inner commutators into currents, $[n_e(\vect{q})^\dagger, H] = \vect{q}\cdot\vect{j}_e(\vect{q})^\dagger$ and $[H, n_\kappa(\vect{q})] = \vect{q}\cdot\vect{j}_\kappa(\vect{q})$. Third, the commutator of $H$ with the nuclear current contains the kinetic energy of the nuclei and the gradient of the potential energy at each nucleus, $\ii M_\kappa^{-1}\sum_{I\in\kappa}\ee^{\ii\vect{q}\cdot\vect{X}_I}\,\vect{q}\cdot\nabla_IU$, with $U$ the potential energy. The electron current commutes with every nuclear operator. It fails to commute only with the electron positions inside the force, and only the electron--nucleus Coulomb energy $\sum_{j,I}w_\kappa(\vect{r}_j - \vect{X}_I)$, $w_\kappa(\vect{y}) = -Z_\kappa\alpha/y$, depends on both. The electron--electron and nucleus--nucleus interactions, and every kinetic term, drop out. What is left involves one function, the mean electron density $\nd_e^{(\kappa)}(\vect{y})$ at the displacement $\vect{y}$ from a nucleus of species $\kappa$, averaged over the nuclei of that species and over their motion. Its value at $\vect{y} = 0$ is the contact density $\nd_e^{(\kappa)}$, the electron density at the nucleus, and far from the nucleus it averages to the mean density $\nd_e$.
\begin{keyresult}[Mixed electron--nucleus third moment]
For point nuclei, nonrelativistic electrons and local potentials, at any $\vect{q}$, temperature and coupling strength,
\begin{equation}
  \mathscr{M}^3_{e\kappa}(\vect{q}) = -\frac{\nd_\kappa}{m_eM_\kappa}\int\dd^3y\;\nd_e^{(\kappa)}(\vect{y})\,\ee^{-\ii\vect{q}\cdot\vect{y}}\,(\vect{q}\cdot\nabla)^2w_\kappa(\vect{y}) .
  \label{eq:8.3-M3}
\end{equation}
At sites of cubic symmetry, for $q$ small compared with the shortest reciprocal-lattice vector $G_{\mathrm{min}}$ but large compared with the inverse size of the sample,
\begin{equation}
  \mathscr{M}^3_{e\kappa}(\vect{q}) = -\frac{Z_\kappa e^2q^2\,\nd_\kappa}{3\,m_eM_\kappa}\bigl[\nd_e^{(\kappa)} + 2\nd_e\bigr] .
  \label{eq:8.3-M3lim}
\end{equation}
The high-frequency interference of electrons and nuclei is negative, and its weight is fixed by the electron density at the nuclei.
\end{keyresult}
\noindent The third moment of a single species is classical, and Puff wrote it for an interacting Bose fluid~\cite{puff1965application}. The two terms of the bracket have different origins, which the identity
\begin{equation*}
  (\vect{q}\cdot\nabla)^2\frac{1}{y} = -\frac{4\pi}{3}\,q^2\,\delta^3(\vect{y}) + q^2\,\frac{3(\hat{\vect{q}}\cdot\hat{\vect{y}})^2 - 1}{y^3}
\end{equation*}
makes visible. The contact part picks out the density at the nucleus and gives $\nd_e^{(\kappa)}/3$ at every $q$, with the phase $\ee^{-\ii\vect{q}\cdot\vect{y}}$ equal to one on its support. The quadrupolar part averages to zero against any density with cubic symmetry about the nucleus that varies on scales short compared with $1/q$, the core and the valence bonds alike. Against the uniform part of the density, however, it does not vanish. That integral is dominated by distances $y \sim 1/q$, where the phase matters, and since the kernel scales as $1/y^3$ it is independent of $q$; it contributes $2\nd_e/3$. It is the Coulomb field of the bare nucleus acting on the electrons of the whole crystal at distances $\sim 1/q$. The order of limits is the physical one: the sample is infinite at fixed $q$, and only then is $q$ small.

The sign has a simple reading. For an electron bound to a nucleus by a spring of constant $K$, \cref{eq:8.3-M3} gives $\mathscr{M}^3_{e\kappa} = -Kq^2/(m_eM_\kappa \Vol)$ at small $q$. The relative coordinate makes the sign obvious: in the mode in which the two particles oscillate against each other, at $\omega^2 = K(m_e^{-1} + M_\kappa^{-1})$, a displacement of one is accompanied by an opposite displacement of the other (\cref{app:D.26}). The Coulomb attraction is a spring whose stiffness, the curvature $\nabla^2w_\kappa = Z_\kappa e^2\delta^3(\vect{y})$, sits at the nucleus, and the contact density measures how much of the electron cloud feels it. Together with the vanishing first moment, the negative third moment says that the interference is positive at low energies, where the electrons follow the nucleus (\cref{sec:8.6}), and negative at high energies, where they are left behind. The positive part sits in the phonon final states, on the acoustic branch, at energies that vanish with $q$, and on the optical phonons, wherever the electrons move with the nuclei (\cref{sec:8.6}). Every electronic final state carries the negative interference $-(m_e/M_\kappa)|\langle f|n_e|0\rangle|^2$ at small $q$, as \cref{eq:8.3-boostward} below shows.

\Cref{eq:8.3-M3lim} has a limited range in $q$, and three regimes follow from the exact form, which \cref{app:D.26} writes as a sum over reciprocal-lattice vectors. For $q \ll G_{\mathrm{min}}$, \cref{eq:8.3-M3lim} holds. Its real part is corrected at relative order $(q/G_{\mathrm{min}})^2$ at every cubic site. At sites without inversion symmetry, such as the tetrahedral sites of silicon, there is also a correction of first order in $q/G_{\mathrm{min}}$. It is purely imaginary. It comes from the part of the density odd about the nucleus, whose Fourier components are imaginary, and it enters $\Imag\mathscr{M}^3_{e\kappa}$ alone. In silicon $G_{\mathrm{min}} = |\vect{G}_{111}| = \SI{3.95}{keV}$ (\cref{sec:3.1}). Once $q$ passes $G_{\mathrm{min}}$ the valence structure starts to matter, and the bracket depends on $q$, although the contact term still dominates. Once $q$ reaches the inverse size $1/a_{\mathrm{core}}$ of the innermost shell, the core itself is resolved: the quadrupolar kernel then acts on the density near the nucleus as on a uniform one, and the bracket climbs from $\nd_e^{(\kappa)}/3$ toward $\nd_e^{(\kappa)}$. For a density with spherical symmetry about the nucleus, which is exact for a free atom and a good approximation for the core in a cubic crystal, the angular integral of \cref{eq:8.3-M3} can be done once and for all (\cref{app:D.26}):
\begin{equation}
  \mathscr{M}^3_{e\kappa}(\vect{q}) = -\frac{Z_\kappa e^2q^2\,\nd_\kappa}{m_eM_\kappa}\Bigl[\frac{\nd_e^{(\kappa)}}{3} + 2\int_0^\infty\dd y\;\frac{\bar\nd_e^{(\kappa)}(y)\,j_2(qy)}{y}\Bigr] ,
  \label{eq:8.3-M3sph}
\end{equation}
with $\bar\nd_e^{(\kappa)}(y)$ the spherically averaged density at distance $y$ from the nucleus and $j_2$ the spherical Bessel function of order two. Since $\int_0^\infty\dd x\,j_2(x)/x = 1/3$, a density that is uniform over the range $y \sim 1/q$ contributes $2/3$ of its value through the integral. This is the origin of the $2\nd_e/3$ of the small-$q$ limit.

The contact density is a material quantity in its own right, and it is computed from all-electron densities. Pseudopotential calculations replace the cores and cannot supply it; the projector-augmented-wave reconstruction or a full-potential code restores the density near the nucleus. The spherical atomic part comes from a calculation of the free atom, and the crystal adds a small anisotropy of the valence (\cref{app:E.17}). Relativity changes the density at the nucleus more than anything else. For a finite nucleus it raises the contact density by about \SI{8}{\percent} in silicon, an effect that grows roughly as $(Z\alpha)^2$ in heavier atoms and that \cref{sec:12.5} adds to \cref{eq:8.3-M3lim}.

\begin{workedexample}{The mixed third moment of silicon}\label{we:8-third}
We evaluate \cref{eq:8.3-M3sph} for silicon at $q = \SI{10}{keV}$, a typical momentum transfer for the Migdal effect. For the density we take the free atom in its ground state from an all-electron Hartree--Fock calculation, nonrelativistic and with a point nucleus. Its total energy, $-288.854$~hartree, is the Roothaan--Hartree--Fock value quoted in \cref{sec:4.3}~\cite{bunge1993roothaan}. The contact density is $\nd_e^{(\kappa)} = 1765\,a_0^{-3} = \SI{1.19e28}{cm^{-3}}$, converged to \SI{0.1}{\percent} in the Gaussian basis, of which the ten core electrons contribute $1757.7\,a_0^{-3}$ and the four valence electrons, through their 3s orbital, $7.6\,a_0^{-3}$, or \SI{0.43}{\percent}. Against it, the mean density of all fourteen electrons, \SI{7.00e23}{cm^{-3}}, is tiny: $2\nd_e/\nd_e^{(\kappa)} = 1.2\times10^{-4}$. In silicon the term that comes from the uniform density is a correction at the level of $10^{-4}$.

At $q = \SI{10}{keV} = 2.68\,a_0^{-1}$ the kernel $j_2(qy)/y$ samples distances $y \sim 1/q = \SI{0.20}{\angstrom}$. This is above $G_{\mathrm{min}} = \SI{3.95}{keV}$, in the middle regime, and below the inverse size $\sim 14/a_0 = \SI{52}{keV}$ of the 1s shell. The bracket of \cref{eq:8.3-M3sph} comes out at $0.3353\,\nd_e^{(\kappa)}$, against $1/3$ in the limit $q \to 0$: \cref{eq:8.3-M3lim} is \SI{0.6}{\percent} low at \SI{10}{keV}. The core supplies \SI{99.56}{\percent} of the bracket and the valence \SI{0.44}{\percent}, almost all of it the 3s density at the nucleus. The crystal valence, which the free atom misses, enters at most through the mean density, $1.8\times10^{-4}$ of the bracket. As $q$ grows the 1s shell is resolved; the bracket is $0.347\,\nd_e^{(\kappa)}$ at \SI{30}{keV} and $0.426\,\nd_e^{(\kappa)}$ at \SI{100}{keV}.

The mixed moment is best read against the nuclear f-sum $\mathscr{M}^1_{\kappa\kappa} = q^2\nd_\kappa/M_\kappa$. Their ratio is a squared energy,
\begin{equation*}
  -\frac{\mathscr{M}^3_{e\kappa}}{\mathscr{M}^1_{\kappa\kappa}} = \frac{Z_\kappa e^2}{m_e}\times(\text{bracket}) = (\SI{8.78}{keV})^2 ,
\end{equation*}
with $(\SI{8.76}{keV})^2$ in the limit $q \to 0$. The scale is that of the K shell: the 1s orbital energy of the same calculation is \SI{1.87}{keV}, and the 1s electrons move with kinetic energies near \SI{2.7}{keV} (\cref{sec:4.3}). The third moment of the electron--nucleus interference is a property of the core electrons of silicon, to better than a percent, and a calculation that freezes the core or replaces it by a pseudopotential gets it wrong by construction.
\end{workedexample}

The diagonal third moments follow from the same algebra. For species $s$ of mass $m_s$,
\begin{equation}
  \mathscr{M}^3_{ss}(\vect{q}) = \frac{\nd_sq^2}{m_s}\Bigl[\Bigl(\frac{q^2}{2m_s}\Bigr)^2 + \frac{3q^2\langle(\hat{\vect{q}}\cdot\vect{p})^2\rangle_s}{m_s^2}\Bigr]
  + \frac{1}{m_s^2\Vol}\Bigl\langle\sum_{i,j\in s}\ee^{-\ii\vect{q}\cdot(\vect{r}_i - \vect{r}_j)}(\vect{q}\cdot\nabla_i)(\vect{q}\cdot\nabla_j)\,U\Bigr\rangle ,
  \label{eq:8.3-M3ss}
\end{equation}
with $\langle(\hat{\vect{q}}\cdot\vect{p})^2\rangle_s$ the mean square of one particle's momentum along $\hat{\vect{q}}$, $2m_s\bar K/3$ for electrons in a cubic crystal, with the $\bar K$ of \cref{eq:4.3-mv} (\cref{app:D.26}). The kinetic bracket is the free recoil, broadened by the motion of the particles. The potential term measures the curvature of the energy surface on which the particles move, and for electrons in jellium it reduces to $\nd_eq^2\omega_p^2/m_e$: the third moment of the electron gas is the plasma frequency squared times the f-sum. In a crystal with mobile nuclei, however, the electron--nucleus curvature at the nuclei dominates it, and the next subsection shows that it is tied to \cref{eq:8.3-M3lim} exactly.

\subsection{Translation Ward identities}\label{sec:8.3-ward}

A dark-matter coupling that is the same per unit mass for every particle, a heavy scalar coupled to nucleons in a crystal of one element for instance, pushes every part of the crystal in the same proportion. At long wavelength the push is uniform, like a gravitational field, and by the equivalence principle it accelerates every particle equally: it can move the crystal, or at finite $q$ launch a sound wave, but it cannot excite the crystal's internal motion. The identities that make this precise follow from translation invariance, which in quantum mechanics is the conservation of total momentum.

Add the continuity equations of all species, each weighted by its mass. The mass-density operator $n_{\mathrm{M}}(\vect{q}) = \sum_sm_sn_s(\vect{q})$, with $s$ running over the electrons and every nuclear species, then obeys
\begin{equation}
  [H, n_{\mathrm{M}}(\vect{q})] = \vect{q}\cdot\vect{g}(\vect{q}) , \qquad
  \vect{g}(\vect{q}) = \sum_sm_s\vect{j}_s(\vect{q}) = \sum_i\frac12\bigl\{\vect{p}_i,\,\ee^{\ii\vect{q}\cdot\vect{r}_i}\bigr\} ,
  \label{eq:8.3-mass}
\end{equation}
where the last sum runs over all particles, electrons and nuclei, with their positions and momenta. The operator $\vect{g}(\vect{q})$ is the momentum density, and at $\vect{q} = 0$ it is the total momentum. Between eigenstates, $[H, X]$ becomes $\omega_{fi}X$, so $\omega_{fi}\langle f|n_{\mathrm{M}}(\vect{q})|i\rangle = \vect{q}\cdot\langle f|\vect{g}(\vect{q})|i\rangle$, and inserted into \cref{eq:4.2-Sab} this is a statement about every response that involves the mass density.
\begin{keyresult}[Translation Ward identity]
For any operator $O_B$ and $\omega \neq 0$, with $S_{Bs}$ the response matrix of $O_B$ and $n_s$,
\begin{equation}
  \sum_sm_s\,S_{Bs}(\vect{q},\omega) = \frac{\vect{q}}{\omega}\cdot S_{B\vect{g}}(\vect{q},\omega) .
  \label{eq:8.3-ward}
\end{equation}
In amplitudes, $\omega_{fi}\sum_sm_s\langle f|n_s(\vect{q})|i\rangle = \vect{q}\cdot\langle f|\vect{g}(\vect{q})|i\rangle$ for every pair of states.
\end{keyresult}
\noindent The name follows field theory, where a Ward identity is a conservation law written for matrix elements. Those of quantum electrodynamics, and the density identity \cref{eq:7.2-ward} of the superconductor, express the conservation of charge; the identities here express the conservation of momentum. \Cref{eq:8.3-ward} holds under the hypothesis of the f-sum rule, that every interaction depends on positions and spins only.

So far the identity only adds up continuity equations. Its content comes from the next commutator. The total momentum $\vect{g}(\vect{0})$ commutes with $H$, since the crystal holds itself together and no external potential acts. For interactions between pairs of particles, the Coulomb interactions among them, $[H, \vect{g}(\vect{q})]$ is then not merely small at small $q$ but proportional to $\vect{q}$ (\cref{app:D.26}):
\begin{equation}
  [H, g^\alpha(\vect{q})] = q_\beta\,\mathsf{T}^{\alpha\beta}(\vect{q}) ,
  \label{eq:8.3-flux}
\end{equation}
with $\mathsf{T}^{\alpha\beta}$ the momentum-flux tensor, the kinetic flux of momentum plus the stress carried by the forces between pairs. This is Newton's third law in operator form: the forces inside the crystal add up to zero, and momentum can only flow from one region to the next. Combined with \cref{eq:8.3-ward}, it gives each factor of $\vect{q}/\omega$ a partner.
\begin{keyresult}[Soft acoustic limit]
For pairwise, translation-invariant interactions and $\omega \neq 0$,
\begin{equation}
  \sum_sm_s\,S_{Bs}(\vect{q},\omega) = \frac{q_\alpha q_\beta}{\omega^2}\,S_{B\mathsf{T}^{\alpha\beta}}(\vect{q},\omega) , \qquad
  S_{n_{\mathrm{M}}n_{\mathrm{M}}}(\vect{q},\omega) = \frac{q_\alpha q_\beta q_\gamma q_\delta}{\omega^4}\,S_{\mathsf{T}^{\alpha\beta}\mathsf{T}^{\gamma\delta}}(\vect{q},\omega) .
  \label{eq:8.3-soft}
\end{equation}
At fixed $\omega \neq 0$ the response of the mass density vanishes as $q^4$, two powers faster than that of any single species. As $q \to 0$ its weight collects on the acoustic branches, where $\omega = c_\nu q$ vanishes with $q$.
\end{keyresult}
\noindent On the acoustic branch $q^4/\omega^4 = c_\nu^{-4}$ stays finite, and the mass density makes sound. Everywhere else, at an optical phonon, a pair of phonons with a finite total energy, or an electronic excitation, its amplitude is $O(q^2)$, against $O(q)$ for the density of one species. For a coupling proportional to mass, the single optical phonon, the two-phonon final states and the anharmonic channel therefore have structure factors $\propto q^4$, and the single acoustic phonon $\propto q$~\cite[Table~I and Sect.~V\,B]{campbelldeem2020multiphonon}, anharmonic terms included. In the harmonic lattice it is the acoustic sum rule \cref{eq:5.2-asr}: a rigid displacement exerts no force, which is the statement that the potential part of $[H, \vect{g}(\vect{q})]$ vanishes at $\vect{q} = 0$.

The identity sorts any coupling to the densities. Write $\sum_sc_sn_s = \lambda n_{\mathrm{M}} + \sum_s(c_s - \lambda m_s)n_s$ for any $\lambda$. The first term reaches internal excitations at order $q^2$ alone. At long wavelengths, as a result, only the part of a coupling that is not proportional to mass shakes the inside of the crystal. A coupling to nucleon number differs from one to mass only by the electron masses and the nuclear binding energies, and for it the optical phonons of a polar crystal are suppressed accordingly at small $q$.

The electron masses give the most instructive case. In a crystal of one element, the density of all nuclei is $\sum_\kappa n_\kappa = (n_{\mathrm{M}} - m_en_e)/M_\kappa$. For every final state $f$ at a finite energy $\omega_{f0}$ above the ground state, \cref{eq:8.3-soft} then makes the mass density drop out to the order that matters:
\begin{equation}
  \Bigl\langle f\Bigl|\sum_\kappa n_\kappa(\vect{q})\Bigr|0\Bigr\rangle = -\frac{m_e}{M_\kappa}\,\langle f|n_e(\vect{q})|0\rangle + O(q^2) .
  \label{eq:8.3-boostward}
\end{equation}
A nucleus that is kicked excites the electrons as if they had been kicked with the momentum $m_e\vect{q}/M_\kappa$, in the opposite direction. This is the boost of the Migdal effect, \cref{eq:8.6-boost} summed over all nuclei, now as an exact consequence of momentum conservation (\cref{app:D.26}). At small $q$ the electronic part of the nuclear response is $(m_e/M_\kappa)^2S_{ee}$, $3.8\times10^{-10}$ times the electron structure factor in silicon, and the interference $S_{e\kappa}$ summed over the nuclei is $-(m_e/M_\kappa)S_{ee}$ there: negative, as the third moment requires.

The third moments obey the same identity. Multiplying \cref{eq:8.3-ward} by $\omega^3$ and integrating leaves one commutator with $n_e^\dagger$ acting on $q_\alpha q_\beta\mathsf{T}^{\alpha\beta}$, so $m_e\mathscr{M}^3_{ee} + \sum_\kappa M_\kappa\mathscr{M}^3_{e\kappa}$ is $O(q^3)$ at least. The kinetic part of \cref{eq:8.3-M3ss} is of order $q^4$, and the terms of order $q^2$ must cancel between the electron self moment and the mixed moments. They do (\cref{app:D.26}). The potential term of \cref{eq:8.3-M3ss} for electrons is, at small $q$ and cubic sites,
\begin{equation}
  \mathscr{M}^3_{ee}(\vect{q}) \to \frac{e^2q^2}{3m_e^2}\Bigl[\sum_\kappa Z_\kappa\nd_\kappa\nd_e^{(\kappa)} + 2\nd_e^2\Bigr] ,
  \label{eq:8.3-M3ee}
\end{equation}
which by charge neutrality, $\sum_\kappa Z_\kappa\nd_\kappa = \nd_e$, is exactly $-\sum_\kappa(M_\kappa/m_e)$ times \cref{eq:8.3-M3lim}. The electron self moment carries the same contact densities. In jellium, where the nuclei are replaced by a fixed uniform background and translation invariance is lost, the contact density is replaced by $\nd_e$, and the bracket $3\nd_e^2$ returns Puff's $\nd_eq^2\omega_p^2/m_e$. For silicon at small $q$ the potential term divided by the f-sum is $e^2[\nd_e^{(\kappa)} + 2\nd_e]/3m_e = (\SI{2.34}{keV})^2$, far above the kinetic terms $(q^2/2m_e)^2 = (\SI{98}{eV})^2$ and $2q^2\bar K/m_e = (\SI{469}{eV})^2$ at $q = \SI{10}{keV}$.

\begin{workedexample}{The Ward identity of a harmonic chain}\label{we:8-chain}
Take a chain of identical nuclei of mass $M$, joined by springs of constant $K$ at spacing $d$, with positions $x_l = ld + u_l$, momenta $P_l$ and periodic boundary conditions over $N$ sites. Its single branch is $\omega_k = 2(K/M)^{1/2}|\sin(kd/2)|$. For numbers we use the longitudinal $[100]$ chain of silicon planes of \cref{sec:5.2}: $M = \SI{26.16}{GeV}$, $d = a/4 = \SI{1.358}{\angstrom}$ and a maximum frequency $2(K/M)^{1/2} = \SI{64.1}{meV}$, which give the sound speed $c = d(K/M)^{1/2} = \SI{6.61}{km.s^{-1}}$. The two operators of \cref{eq:8.3-ward} are $n(q) = \sum_l\ee^{\ii qx_l}$ and $g(q) = \sum_l\frac12\{P_l, \ee^{\ii qx_l}\}$, and the identity reads $M\omega_{fi}\langle f|n|i\rangle = q\langle f|g|i\rangle$.

\emph{One phonon.} With $u_l = \sum_k(2NM\omega_k)^{-1/2}(a_k\ee^{\ii kld} + a_k^\dagger\ee^{-\ii kld})$, the term of first order in $u_l$ creates the phonon $k = q$ from the ground state with $\langle1_q|n|0\rangle = \ii q\,(N/2M\omega_q)^{1/2}$. The momentum density to the same order is $\sum_l\ee^{\ii qld}P_l$, and $\langle1_q|P_l|0\rangle = \ii M\omega_q\langle1_q|u_l|0\rangle$ gives $\langle1_q|g|0\rangle = \ii(NM\omega_q/2)^{1/2}$. The ratio is $q/M\omega_q$, as the identity demands. At $q = \SI{1}{keV} = \SI{0.507}{\angstrom^{-1}}$, $qd = 0.688$ and $\omega_q = \SI{21.6}{meV}$.

\emph{Two phonons.} The second-order terms are $-\frac12q^2\sum_l\ee^{\ii qld}u_l^2$ in $n$ and $\frac12\ii q\sum_l\ee^{\ii qld}\{P_l, u_l\}$ in $g$. For the pair $|k, q - k\rangle$, with $k \ne q - k$, they give
\begin{equation*}
  \langle k, q - k|n|0\rangle = -\frac{q^2}{2M(\omega_k\omega_{q-k})^{1/2}} , \qquad
  \langle k, q - k|g|0\rangle = -\frac{q\,(\omega_k + \omega_{q-k})}{2(\omega_k\omega_{q-k})^{1/2}} ,
\end{equation*}
and $M(\omega_k + \omega_{q-k})$ times the first is $q$ times the second. The pair has a finite energy as $q \to 0$, near $2\omega_k$, and its amplitude is of order $q^2$: the soft acoustic limit, here visible term by term. Measured against the one-phonon weight, the two-phonon weight carries a factor of order $q^2/(2M\omega)$, the recoil energy over a phonon energy, $8.8\times10^{-4}$ at \SI{1}{keV}.

\emph{The flux.} The forces enter through $[H, g(q)]$. Each spring pulls its two ends with opposite forces, $\pm K\Delta_l$ with $\Delta_l = u_{l+1} - u_l$, and the potential part of the commutator is $\ii K\sum_l\Delta_l(\ee^{\ii qx_{l+1}} - \ee^{\ii qx_l})$. The difference of the phases, $\ee^{\ii qld}(\ee^{\ii qd} - 1)$ at lowest order, carries the factor $q$ of \cref{eq:8.3-flux}: Newton's third law makes the force on the chain as a whole vanish, the one-dimensional form of \cref{eq:5.2-asr}. Its one-phonon element is $4\ii K\sin^2(qd/2)\,(N/2M\omega_q)^{1/2}$, and with $\omega_q\langle1_q|g|0\rangle = \langle1_q|[H, g]|0\rangle$ we recover $M\omega_q^2 = 4K\sin^2(qd/2)$, the dispersion of the chain, up to terms of relative order $q^2/(M\omega_q)$ from the kinetic part of the flux. The Ward identity, read at one phonon, is the equation of motion.

A cubic term in the springs changes none of this. It adds to the two-phonon amplitude a path through the one-phonon state $|1_q\rangle$, whose amplitude from $n$ is of order $q$. The anharmonic path alone is of order $q$ times the cubic coupling. The flux \cref{eq:8.3-flux} still carries its factor $q$, because the cubic springs are still forces between pairs, and the total amplitude to the pair therefore remains of order $q^2$: the terms of order $q$ cancel between the paths.
\end{workedexample}

\subsection{Table of exact relations}\label{sec:8.3-table}

\Cref{tab:8-exact} collects the exact relations of \cref{sec:4.2,sec:4.3} and of this section with their hypotheses. Each follows from stationarity, the positivity of probabilities, causality, or the conservation of particle number or momentum, and together they constrain the nuisance parameters with which \cref{sec:8.8} turns truncations into uncertainties.

\begin{table}[tbp]
  \centering
  \small
  \caption{Relations that hold for every target, with the hypotheses they need. ``Stationary'' means any state diagonal in the energy basis, the Gibbs state at any temperature included; ``local'' means that every interaction depends on positions and spins only; ``translation invariant'' means pairwise interactions, no external potential.}
  \label{tab:8-exact}
  \begin{tabular}{@{}p{0.22\linewidth}p{0.44\linewidth}p{0.26\linewidth}@{}}
    \toprule
    Relation & Statement & Hypotheses \\
    \midrule
    Detailed balance & $S_{ab}(-\vect{q},-\omega) = \ee^{-\beta\omega}S_{ba}(\vect{q},\omega)$, \cref{eq:A.7-DB} & Gibbs state \\
    Fluctuation--dissipation & $S$ from $\Imag R$, \cref{eq:4.2-fdt} & Gibbs state \\
    Kramers--Kronig & $\Real R$ from $\Imag R$, \cref{eq:4.3-KK} & causality \\
    Positivity & $S_{ab}$ positive semidefinite, \cref{eq:4.2-positivity} & stationary \\
    Interference bound & $|\varGamma_{AB}| \le 2(\varGamma_A\varGamma_B)^{1/2}$, \cref{eq:8.3-bound} & stationary \\
    Species f-sum & $\mathscr{M}^1_{ss'} = \delta_{ss'}q^2\nd_s/m_s$, \cref{eq:4.3-species} & stationary, local \\
    Moment matrices & \cref{eq:8.3-moments}, Hankel positivity \cref{eq:8.3-hankel} & stationary \\
    Mixed third moment & $\mathscr{M}^3_{e\kappa}$, \cref{eq:8.3-M3,eq:8.3-M3lim} & stationary, local, point nuclei \\
    Ward identity & $\sum_sm_sS_{Bs} = (\vect{q}/\omega)\cdot S_{B\vect{g}}$, \cref{eq:8.3-ward} & stationary, local \\
    Soft acoustic limit & $S_{n_{\mathrm{M}}n_{\mathrm{M}}} = O(q^4)$ at $\omega \ne 0$, \cref{eq:8.3-soft} & translation invariant \\
    \bottomrule
  \end{tabular}
\end{table}

None of these relations involves the coupling to the dark matter beyond the positivity of $\mathfrak{D}_{ab}$, and none involves screening.
